\begin{proof}[Proof of \cref{obj:lem:identification-infrastructure}]
\begin{enumerate}
\item For each \(a\in\{0,1\}\), consistency in
\cref{obj:ass:consistency} gives
\[
P(A=a,Y=1)=P(A=a,Y(a)=1).
\]
To calculate the latter probability, define the potential-coordinate vector and its relevant finite subset by
\[
\mathbf V_{\mathrm{pot}}
:=(X,S(0),S(1),Y(0),Y(1)),
\qquad
\mathcal V_{\mathrm{pot},a}
:=
\left\{
(x,s_0,s_1,y_0,y_1)
\in\{0,\ldots,d-1\}\times\{0,1\}^{4}
:\ y_a=1
\right\},
\quad a\in\{0,1\}.
\]
Summing the atomwise factorization supplied by
\cref{obj:ass:randomized-independence} over this subset yields
\[
\begin{aligned}
P(A=a,Y(a)=1)
&=\sum_{\boldsymbol v\in\mathcal V_{\mathrm{pot},a}}
  P(A=a,\mathbf V_{\mathrm{pot}}=\boldsymbol v)\\
&=P(A=a)
  \sum_{\boldsymbol v\in\mathcal V_{\mathrm{pot},a}}
  P(\mathbf V_{\mathrm{pot}}=\boldsymbol v)\\
&=P(A=a)P(Y(a)=1).
\end{aligned}
\]
Here the sums partition the indicated events by their potential-coordinate tuples.
By \cref{obj:ass:balanced-randomization} and the binary treatment partition,
\[
P(A=1)=\frac12,
\qquad
P(A=0)=1-P(A=1)=\frac12.
\]
Consequently, the two arm identities are
\[
2P(A=1,Y=1)=P(Y(1)=1),
\qquad
2P(A=0,Y=1)=P(Y(0)=1).
\]
% lean: CausalSmith.Stat.MarNearcompleteFrontier.randomized_arm_outcome

\item Since both potential outcomes are binary, their numerical values equal their success indicators. Linearity of expectation therefore gives
\[
\begin{aligned}
\tau(P)
&=\mathbb E_P\!\left[
  \mathbf 1_{\{Y(1)=1\}}-\mathbf 1_{\{Y(0)=1\}}
  \right]\\
&=P(Y(1)=1)-P(Y(0)=1).
\end{aligned}
\]
% lean: CausalSmith.Stat.MarNearcompleteFrontier.tau_eq_potential_mass

\item Define the cell events, for their full range of indices, by
\[
\mathcal E_{x,a,s}:=\{X=x,A=a,S=s\},
\qquad
x\in\{0,\ldots,d-1\},\quad a,s\in\{0,1\}.
\]
The observation map preserves \(X,A,S,R\), and on \(\{R=1\}\) its outcome mark equals \(Y\). Thus its image law satisfies
\[
\begin{gathered}
P_O(\mathcal E_{x,a,s})=P(\mathcal E_{x,a,s}),
\qquad
P_O(A=a)=P(A=a),\\
P_O(\mathcal E_{x,a,s},R=1)
  =P(\mathcal E_{x,a,s},R=1),\\
P_O(\mathcal E_{x,a,s},R=1,RY=1)
  =P(\mathcal E_{x,a,s},R=1,Y=1).
\end{gathered}
\]
In particular, \cref{obj:def:identified-functional} supplies
\[
\frac{P_O(\mathcal E_{x,a,s})}{P_O(A=a)}
=\frac{P(\mathcal E_{x,a,s})}{P(A=a)},
\qquad
\mu_P(x,a,s)
=\frac{P(\mathcal E_{x,a,s},R=1,Y=1)}
       {P(\mathcal E_{x,a,s},R=1)},
\]
with zero assigned to every ratio whose denominator is zero.

We identify the weighted contribution by splitting according to the cell mass. If \(P(\mathcal E_{x,a,s})=0\), nonnegativity and event inclusion imply
\[
0\le P(\mathcal E_{x,a,s},Y=1)
\le P(\mathcal E_{x,a,s})=0.
\]
Hence both the weighted observed contribution and
\(P(\mathcal E_{x,a,s},Y=1)/P(A=a)\) are zero.

If \(P(\mathcal E_{x,a,s})>0\), event inclusion gives
\[
0<P(\mathcal E_{x,a,s})\le P(A=a).
\]
The arrival floor in \cref{obj:ass:arrival-floor} and the stipulated range of \(q\) give
\[
\frac{P(\mathcal E_{x,a,s},R=1)}
     {P(\mathcal E_{x,a,s})}
\ge q\ge\frac12>0,
\qquad
P(\mathcal E_{x,a,s},R=1)>0.
\]
Conditional independence in \cref{obj:ass:mar}, applied to \(R=1\) and \(Y=1\), gives the cross-product identity
\[
P(\mathcal E_{x,a,s},R=1,Y=1)\,
P(\mathcal E_{x,a,s})
=
P(\mathcal E_{x,a,s},R=1)\,
P(\mathcal E_{x,a,s},Y=1).
\]
Dividing by the two positive cell masses yields
\[
\mu_P(x,a,s)
=
\frac{P(\mathcal E_{x,a,s},R=1,Y=1)}
     {P(\mathcal E_{x,a,s},R=1)}
=
\frac{P(\mathcal E_{x,a,s},Y=1)}
     {P(\mathcal E_{x,a,s})}.
\]
Multiplication by the cell weight now proves, in both cases,
\[
\frac{P_O(\mathcal E_{x,a,s})}{P_O(A=a)}
\,\mu_P(x,a,s)
=
\frac{P(\mathcal E_{x,a,s},Y=1)}{P(A=a)}.
\]
% lean: CausalSmith.Stat.MarNearcompleteFrontier.cellContribution_identified

\item For either treatment arm, every realization with \(A=a\) and \(Y=1\) belongs to exactly one cell indexed by its values of \(X\) and \(S\). The disjoint finite-cell partition therefore gives
\[
\sum_{x=0}^{d-1}\sum_{s=0}^{1}
P(\mathcal E_{x,a,s},Y=1)
=
P(A=a,Y=1),
\qquad a\in\{0,1\}.
\]
% lean: CausalSmith.Stat.MarNearcompleteFrontier.cellOutcome_sum

\item Substituting the cell-contribution identity into
\cref{obj:def:identified-functional}, distributing the finite sums, and applying the arm partitions gives
\[
\begin{aligned}
\Psi(P_O)
&=\sum_{x=0}^{d-1}\sum_{s=0}^{1}
\left\{
\frac{P(\mathcal E_{x,1,s},Y=1)}{P(A=1)}
-
\frac{P(\mathcal E_{x,0,s},Y=1)}{P(A=0)}
\right\}\\
&=\frac{P(A=1,Y=1)}{P(A=1)}
  -\frac{P(A=0,Y=1)}{P(A=0)}\\
&=2P(A=1,Y=1)-2P(A=0,Y=1)\\
&=P(Y(1)=1)-P(Y(0)=1)\\
&=\tau(P).
\end{aligned}
\]
The null-cell case above establishes the stipulated zero contributions. When \(d=0\), the baseline-label set is empty, so a probability law with \(X\) taking values in that set cannot exist; the assertion for that case is vacuous.
% lean: CausalSmith.Stat.MarNearcompleteFrontier.tau_eq_psi
\end{enumerate}
\end{proof}
